\section{Failure Mode and Effects Analysis (\acs{FMEA})}
\label{sec:fmea}

This section is a living catalogue of identified failure modes, their
effects and the mitigation currently in place -- hardware, firmware, or
both. It is expected to be revised continuously as the design changes: a
row is updated in place once its mitigation changes, and closed items are
kept (marked minor) rather than deleted, so the table also documents what
has already been addressed. Current scope is electrical and firmware
failure modes; mechanical/enclosure failure modes (ingress, mounting,
connector wear, \ldots) are not yet covered and should be added as a
separate subsection once that design matures.

Each failure mode is tagged HW, FW or HW+FW to show
which layer(s) it originates in or is mitigated by. Row colour is the
severity:

\begin{itemize}
    \item \colorbox{lred}{\phantom{xx}} \textbf{Critical} -- risk of injury (electric shock, fire, cell venting), of a safety requirement (\autoref{sec:fmea} cross-references \secref{sec_ovp} etc.) failing silently, or of hardware damage that a user could not reasonably anticipate.
    \item \colorbox{orange!30}{\phantom{xx}} \textbf{Moderate} -- functional or measurement degradation, false readings, or loss of availability, with no direct safety exposure because another layer still backs it up.
    \item \colorbox{lgreen}{\phantom{xx}} \textbf{Minor} -- cosmetic/nuisance impact, or a previously higher-severity item that has since been mitigated and is kept here as a closed record.
\end{itemize}

\FloatBarrier

\subsection{Battery and \acs{BMS}}

\begin{table}[h!]
    \begin{FMEATable}
    \FMEAItemMinor
        {[HW+FW, closed] Cell overvoltage during charge}
        {A cell exceeds its safe upper voltage limit, e.g.\ through a runaway charge-voltage control loop or a charger fault.}
        {Accelerated cell degradation; risk of thermal event in extreme cases.}
        {\ac{BMS} (BQ76905) autonomously enforces per-cell overvoltage protection (own comparators/FET driver, independent of the host microcontroller) opening the charge \ac{FET}; firmware additionally checks BMS cell/stack voltage each tick and exits CHARGE (\texttt{statemachine.c}), see \ReqRef{R:BMS}.}
    \FMEAItemMinor
        {[HW+FW] Pack overcurrent / short circuit}
        {External short across the battery or charging port, or an internal fault driving excess charge/discharge current, see \ReqRef{R:Short}, \ReqRef{R:Overcurrent}.}
        {Cell/wiring damage, fire risk if unmitigated.}
        {\ac{BMS} two-tier discharge-overcurrent, short-circuit-discharge and charge-overcurrent protection, fully autonomous of the host firmware and of an unresponsive microcontroller (\autoref{tab:fmea_mcu_hang}).}
    \FMEAItemCritical
        {[HW+FW] Cell over-/under-temperature}
        {Charging or discharging outside the safe cell temperature window.}
        {Accelerated ageing; safety risk at the extremes.}
        {\ac{BMS} autonomous over-/under-temperature charge and discharge protection thresholds, independent of host firmware, see \ReqRef{R:BMS}. This feature is not tested currently.}
    \FMEAItemModerate
        {[FW] Loss of \acs{I2C} communication to \acs{BMS}}
        {I2C4 bus errors/NACKs (known marginal 2.4\,V level-shifter output, see architecture.md \#3) stop host reads of the \ac{BMS}.}
        {Firmware loses SoC/voltage/current visibility and cannot react to thresholds it reads from the \ac{BMS} (e.g.\ CHARGE exit); display shows stale values.}
        {Queued as \texttt{ERROR\_IIC\_*} (\texttt{error.c}), visible on the \ac{CLI} only -- not persisted or shown on-screen (\texttt{error\_Write()} is an unimplemented stub). \ac{BMS}-internal protection keeps acting as the safety backstop regardless.}
    \end{FMEATable}
    \caption{FMEA -- Battery and BMS}
\end{table}

\FloatBarrier

\subsection{Charging}

\begin{table}[h!]
    \begin{FMEATable}
    \FMEAItemCritical
        {[FW] Untuned PI gains cause control-loop instability}
        {All eight PI loops (\texttt{ctrl\_pid\_table[]} in \texttt{ctrl\_main.c}), including the charge-current/voltage and high-voltage flyback loops, still run placeholder gains -- no bench tuning has been performed.}
        {Overshoot or oscillation of the regulated voltage/current beyond the commanded set-point, potentially exceeding component or cell ratings transiently.}
        {Hardware backstops only at present: \acs{BMS} autonomous protection, \acs{OVP} crowbar (\secref{sec_ovp}). Treat every first activation of a loop that has not yet been bench-tuned as bring-up, not validated behaviour.}
    \FMEAItemModerate
        {[FW] CV-tail regulation target can diverge from the latch/termination sensor}
        {The CC$\to$CV latch and cycle termination both key off the \ac{BMS} StackVoltage against the same threshold (consistent). Once in CV, though, the loop itself regulates towards $4 \times$ the highest individual cell voltage rather than StackVoltage, to protect the most-charged cell under imbalance.}
        {Under significant cell imbalance, the quantity being latched/terminated on (real StackVoltage) and the quantity the CV loop drives towards (extrapolated $4\times$max-cell) can disagree -- untested at the time of writing.}
        {Termination additionally requires the cell-voltage spread to be $\leq$\qty{50}{\milli\volt}, which should catch most of this in practice, but the underlying two-target tension is unresolved. Open item -- watch during bench validation with an intentionally imbalanced pack.}
    \FMEAItemMinor
        {[FW, closed] Coulomb counter / passed-charge unit error}
        {Passed-charge integration used the wrong unit (0.25 userAsec), skewing reported \acs{SoC}/passed charge.}
        {Displayed \ac{SoC}/charge state could mislead the operator about battery status.}
        {Fixed in firmware (commits \texttt{1eb7f9b}, \texttt{1f3783b}). Kept here, marked closed, as the working example of how this table tracks a resolved finding rather than deleting it.}
    \FMEAItemModerate
        {[HW+FW] Charge-start inrush across the output relay}
        {The output relay (K1, normally closed) used to close onto a running converter whose output node was not at the charger voltage, and the charge-current loop started near the maximum-current duty.}
        {Current pulse into the converter output capacitors and the battery at the start of every charge; trips the \ac{BMS} charge-overcurrent protection, stresses the relay contacts (a welded relay has been seen on one unit).}
        {Multi-phase start (\texttt{charge\_seq.c}): relay open and converter run as an open-loop boost to 99\,\% of the charger voltage, relay closed with the SEK duty stepped to the zero-current duty, current reference ramped from \qty{0}{\ampere} over \qty{5}{\second}. The peak current after the close command is printed for bench tuning. Validated on hardware.}
    \FMEAItemModerate
        {[HW+FW] Converter output node not measured during precharge}
        {\texttt{OUT\_LV}, behind the relay, has no \ac{ADC} channel: \texttt{V\_OUT}, \texttt{V\_TERM} and \texttt{I\_OUT} are all on the jack side. The precharge duty is computed open-loop from \texttt{V\_IN}.}
        {Model error (V\_IN accuracy, dead time, unloaded-boost assumption) leaves \texttt{OUT\_LV} off target when the relay closes: residual inrush if too low, reverse current into the charger if too high.}
        {The target is 99\,\% rather than 100\,\% of the charger voltage, all ramps are slow, and the ratio and relay settle time are tunable constants. Recommendation: add an \texttt{OUT\_LV} divider on a spare \ac{ADC} channel to close the loop.}
    \FMEAItemModerate
        {[HW] Output relay stuck open or slow to close}
        {K1 does not connect the converter to the charger after the close command.}
        {No charge current flows; the charge-current PI integrator would wind up and the loop surge when the path finally closes.}
        {Detected in the current ramp: reference above \qty{300}{\milli\ampere} for \qty{2}{\second} with less than \qty{20}{\milli\ampere} flowing aborts the charge and locks out re-entry until the charger is removed.}
    \FMEAItemModerate
        {[HW] Output relay stuck closed (welded)}
        {K1 is closed during the precharge.}
        {The open-loop boost then runs against the connected charger; the precharge duty is below the charger voltage, so current flows into the battery.}
        {Not detected in the precharge phase today; the \ac{BMS} charge-overcurrent protection is the backstop. Recommendation: abort if charge current is measured while the precharge is running.}
    \FMEAItemModerate
        {[HW+FW] Normal CHARGE exit would cut the device's own power during a deep-discharge recovery}
        {A pack discharged far enough that the \ac{BMS} latches a cell-undervoltage (CUV) fault opens its DSG \ac{FET} (Q5); the device is then powered only through K1 (normally closed) feeding the charger voltage into \texttt{OUT\_LV}, with no battery-side supply at all. Every normal CHARGE exit (BMS fault, stuck-open abort, \ac{ESC}, boot) routes through \texttt{statemachine\_switchtoIdle()}, whose ordinary IDLE relay configuration opens K1.}
        {Opening K1 in this state removes the device's only power source mid-recovery: the microcontroller browns out/resets, unable to ever begin recovering the pack without external intervention.}
        {\texttt{statemachine\_switchtoIdle()} keeps K1 closed whenever CUV is latched, regardless of why it was called -- one chokepoint covers boot, fault-exit, stuck-open and \ac{ESC} alike. The recovery charge itself is entered directly (\texttt{statemachine\_enter\_charge\_low\_current()}), skipping the open-loop \texttt{PRECHARGE} boost since \texttt{V\_IN} reads $\approx$\qty{0}{\volt} with the DSG \ac{FET} open, and starts \texttt{CC\_RAMP} at a reduced current (\qty{100}{\milli\ampere}, tunable) that reaches the pack through Q6 (the separate charge \ac{FET}, gated directly by the \ac{BMS}, independent of DSG) and Q5's body diode. Validated on hardware.}
    \FMEAItemMinor
        {[FW] CUV masked from the CHARGE fault-exit check during recovery}
        {While the deep-discharge recovery charge is active, the CUV bit specifically is excluded from the fault check that would otherwise abort CHARGE -- it is exactly the condition the run exists to clear.}
        {If this masking were broader than intended, a genuine unrelated fault occurring during the same window could be missed by firmware's own CHARGE-exit logic.}
        {The mask is scoped to exactly the CUV bit (bit 6 of Safety Status A); any other \ac{BMS} safety fault bit still aborts the charge immediately, same as a normal charge. Independently, the \ac{BMS}'s own autonomous protections (overvoltage, overcurrent, short-circuit, thermal) and the \ac{OVP} crowbar remain active throughout regardless of this firmware-level check.}
    \end{FMEATable}
    \caption{FMEA -- Charging}
\end{table}

\FloatBarrier

\subsection{Power conversion}

\begin{table}[h!]
    \begin{FMEATable}
    \FMEAItemMinor
        {[HW, closed] Primary half-bridge shoot-through}
        {Simultaneous conduction of the high-side/low-side \acs{GaN} switches on the primary buck-boost bridge.}
        {Very high current spike; potential destruction of the LMG2100R044 half-bridges.}
        {Hardware dead-time insertion configured on \ac{HRTIM} Timer B (PRIM) and Timer E (SEK), independent of firmware control-loop behaviour. Note: the secondary half-bridge is deliberately driven into shoot-through by firmware for AMPMETER/crowbar operation -- by design, not a fault, see \secref{sec_lv_converter}.}
    \FMEAItemCritical
        {[HW+FW] Output overvoltage}
        {A runaway control loop (see charging-table entry above) drives \texttt{v\_out}/\texttt{v\_hv} above the safe limit.}
        {Damage to downstream circuitry, or an unexpectedly high voltage on user-accessible output terminals.}
        {\ac{OVP} crowbar shorts the converter secondary once the output exceeds threshold, latching until an explicit \texttt{OVP\_RESET} (\secref{sec_ovp}); \texttt{protection.c} independently polls \texttt{OVP\_N} every tick and forces the output off regardless of mode.}
    \FMEAItemModerate
        {[HW+FW] Output overcurrent coverage differs by mode -- needs review}
        {\texttt{protection.c} only evaluates \texttt{OCP\_N} while ISOMETER is active (``it trips spuriously in other modes'', per the code comment); the \texttt{OCP\_N} circuit is the high-voltage flyback's current limiter (\secref{sec_ocp}, \autoref{fig_hv_conv_i_lim}), not a general converter latch. A separate analog limiter exists on the primary/output current path (\autoref{fig_i_ctrl_feedforward}) but its fault behaviour is not yet written up in \secref{sec_i_ctrl}.}
        {In 60V\_OUT/10A\_OUT/RESISTANCE modes, output overcurrent is caught by the untuned \texttt{BoostIoutLimit} software loop and, ultimately, by \ac{BMS} pack-level protection -- not by a dedicated fast local latch as in ISOMETER.}
        {\ac{BMS} pack overcurrent-discharge protection remains the backstop.}
    \end{FMEATable}
    \caption{FMEA -- Power conversion}
\end{table}

\FloatBarrier

\subsection{Isolation measurement (high voltage)}

\begin{table}[h!]
    \begin{FMEATable}
    \FMEAItemCritical
        {[HW] Isolation current channel offset masks real leakage}
        {A \qty{129}{\micro\ampere} offset on the \texttt{I\_ISO} channel (op-amp input leakage, architecture.md \#5).}
        {Can make the measured insulation resistance appear better than reality, undermining \ReqRef{R:Iso_voltage} exactly when the test should catch a real fault -- a false-safe measurement error.}
        {\texttt{I\_ISO} is one of \texttt{calibration.c}'s six per-unit calibratable channels (\texttt{zeroCal}/\texttt{gainCal}); a unit that has not been zero-calibrated still carries the full offset.}
    \FMEAItemCritical
        {[HW] Operator contact with high-voltage output during isolation test}
        {ISOMETER energises the output terminals up to \qty{500}{\volt} for the resistance test.}
        {Electric shock.}
        {Shrouded \qty{4}{\milli\meter} safety jacks (\ReqRef{R:Output}); \texttt{OUT\_SEL\_HV} output routing is only active in ISOMETER (\texttt{mode\_table.c}); colour-changing screen border while ISOMETER is active as an operator-facing indicator (commit \texttt{1ac3eeb}); \ac{GDT}s on the output terminals (\secref{sec_terminal}).}
    \FMEAItemCritical
        {[FW] High-voltage flyback loop instability on first bring-up}
        {The HVVoltage/HVCurrent PI loops share the untuned-gains issue above; they previously defaulted to zero gain (silently disabling ISOMETER) and now run real but unverified gains.}
        {Overshoot beyond the \qty{500}{\volt} rating of downstream components at the point where the consequences are worst.}
        {\ac{OVP} crowbar and \ac{GDT}s as hardware backstops (\secref{sec_ovp}, \secref{sec_terminal}); treat the first ISOMETER activation after any gain change as a supervised bench test.}
    \end{FMEATable}
    \caption{FMEA -- Isolation measurement}
\end{table}

\FloatBarrier

\subsection{Thermal}

\begin{table}[h!]
    \begin{FMEATable}
    \FMEAItemModerate
        {[FW] Power-stage over-temperature}
        {4 external \ac{NTC}s (secondary, transformer, current-sense, primary) monitored by \texttt{protection.c}.}
        {Component degradation/failure if reached and not acted on.}
        {\texttt{protection.c} classifies OK/WARNING/ERROR per channel with hysteresis every tick; ERROR forces the output off from \texttt{statemachine\_step()} regardless of mode.}
    \end{FMEATable}
    \caption{FMEA -- Thermal}
\end{table}

\FloatBarrier

\subsection{Firmware / system-level}
\label{tab:fmea_mcu_hang}

\begin{table}[h!]
    \begin{FMEATable}
    \FMEAItemCritical
        {[FW] No microcontroller watchdog}
        {\texttt{HAL\_IWDG\_MODULE\_ENABLED} and \texttt{HAL\_WWDG\_MODULE\_ENABLED} are both disabled (\texttt{stm32g4xx\_hal\_conf.h}); nothing resets the microcontroller if the main loop hangs.}
        {If firmware hangs while a converter output is enabled, the output stays in whatever state it was in until a hardware-level protection (\ac{OVP}, \ac{BMS}) intervenes or the device is unplugged -- there is no firmware-level fail-safe of its own.}
        {None in firmware today; the \ac{BMS}'s autonomous protection and the \ac{OVP} crowbar/dead-time hardware stay active even with an unresponsive microcontroller, so the pack- and converter-level physical protections are the actual backstop. \textbf{Recommendation}: enable the microcontroller's independent watchdog timer (IWDG).}
    \FMEAItemModerate
        {[FW] EEPROM calibration/config corruption}
        {The header, \texttt{hardware\_data} and \texttt{calibration} blocks each carry a CRC32 (\texttt{config\_store.c}).}
        {On mismatch, \texttt{config\_store\_init()} silently falls back to default (uncalibrated / placeholder-PID) values for that block; the device keeps operating with degraded measurement accuracy and no on-screen indication -- only a \ac{CLI} \texttt{printf}.}
        {CRC32-per-block check with default fallback prevents operating on garbage data. \textbf{Recommendation}: surface a UI-visible indicator when a block has been defaulted.}
    \FMEAItemMinor
        {[FW] Event / \acs{I2C} queue overflow}
        {The 64-entry event queue and the \ac{I2C} transaction queue are fixed-size and drop new entries when full rather than blocking.}
        {A dropped \texttt{EVENT\_SM\_STEP} delays the next state-machine tick by one 20\,ms period (self-correcting); a dropped \ac{I2C} transaction delays that read to the next cycle.}
        {Queued as \texttt{ERROR\_EVENT\_OVF}/\texttt{ERROR\_IIC\_OVF} (\texttt{error.c}), \ac{CLI}-visible only, not persisted (\texttt{error\_Write()} stub). Sustained overflow would indicate a starvation bug worth investigating, not a one-off.}
    \end{FMEATable}
    \caption{FMEA -- Firmware / system-level}
\end{table}
